\section{Análisis}

A máxima velocidad el 95\,\% de las respuestas llega sin el encabezado
\texttt{5A}. El error depende de la separación entre transferencias, medida en
la PC entre el fin de una llamada y el inicio de la siguiente:

\begin{table}[H]
\centering
\begin{tabular}{@{}lrr@{}}
\toprule
Separación & Tramas & Encabezado correcto \\
\midrule
menos de 20\,$\mu$s     & 98\,845 & 4.5\,\% \\
20 a 60\,$\mu$s         & 1\,139  & 53.6\,\% \\
60 a 100\,$\mu$s        & 7       & 57.1\,\% \\
100\,$\mu$s o más       & 7       & 100\,\% \\
\bottomrule
\end{tabular}
\end{table}

La respuesta de cada trama se arma en la interrupción de fin de la trama
anterior, que además reinicia SPI1 y vuelve a programar el DMA. Con
$F_{cy} = 4$\,MHz esa rutina se estimó en unos 100\,$\mu$s; si la siguiente
trama comienza antes de que termine, el primer byte que sale por MISO no es
el encabezado. Con el reloj actual del dsPIC esta prueba no puede pasar: el
intervalo mínimo entre tramas lo fija la duración de la interrupción, que
debe medirse con el analizador lógico en RD10.
